% VI. Експериментални резултати и анализ.
% ВНИМАНИЕ: този раздел е СКЕЛЕТ. В него няма и не бива да се появи нито едно
% число, преди то да е измерено по методиката от Раздел V. Празните клетки са
% умишлени и се намират с: grep -n 'TODO' chapters/06_results.tex
\section{Експериментални резултати и анализ}
\label{sec:results}

\subsection{Коректност спрямо тестовите набори}

Резултатът от проверката за коректност се представя в \tabref{tab:correctness}. Всеки
непреминал случай се разглежда поотделно и се отнася към една от две причини: липсваща
конструкция, която е извън обявения обхват, или грешка в реализацията. Разграничението
е съществено, защото първата причина е проектно решение, а втората е дефект.

\begin{table}[H]
\caption{Резултати от тестовите набори на W3C}
\label{tab:correctness}
\centering
\fontsize{10}{12}\selectfont
\begin{tabular}{@{}L{3.6cm}rrrr@{}}
\toprule
\textbf{Тестов набор} & \textbf{Общо} & \textbf{Преминали} & \textbf{Непреминали} & \textbf{Пропуснати} \\
\midrule
Turtle       & \TODO{n} & \TODO{n} & \TODO{n} & \TODO{n} \\
N-Triples    & \TODO{n} & \TODO{n} & \TODO{n} & \TODO{n} \\
SPARQL 1.1   & \TODO{n} & \TODO{n} & \TODO{n} & \TODO{n} \\
\bottomrule
\end{tabular}
\end{table}

\TODO{списък на непреминалите случаи с причината за всеки от тях}

\subsection{Еднократна цена: зареждане и обем}

Времето за зареждане и обемът на получените индекси се представят в
\tabref{tab:loading}. Величините се отчитат за една и съща входна редица от триплети
и при еднакви условия по методиката от Раздел~\ref{sec:method}.

\begin{table}[H]
\caption{Време за зареждане, обем на индексите и върхова потребена памет}
\label{tab:loading}
\centering
\fontsize{10}{12}\selectfont
\begin{tabular}{@{}L{3.6cm}rrr@{}}
\toprule
\textbf{Реализация} & \textbf{Зареждане, s} & \textbf{Обем, MB} & \textbf{Памет, MB} \\
\midrule
Trident & \TODO{изм.} & \TODO{изм.} & \TODO{изм.} \\
Jena    & \TODO{изм.} & \TODO{изм.} & \TODO{изм.} \\
RDFLib  & \TODO{изм.} & \TODO{изм.} & \TODO{изм.} \\
\bottomrule
\end{tabular}
\end{table}

\subsection{Повтаряща се цена: изпълнение на заявки}

Времената за изпълнение по заявки се представят в \tabref{tab:queries}. Колоната с
броя решения не е допълнителна информация, а условие за валидност: ред, в който
броевете не съвпадат, не се тълкува като сравнение на скорост.

\begin{table}[H]
\caption{Медианно време за изпълнение по заявки}
\label{tab:queries}
\centering
\fontsize{10}{12}\selectfont
\begin{tabular}{@{}L{2.6cm}rrrr@{}}
\toprule
\textbf{Заявка} & \textbf{Trident, ms} & \textbf{Jena, ms} & \textbf{RDFLib, ms} & \textbf{Решения} \\
\midrule
\TODO{заявка} & \TODO{изм.} & \TODO{изм.} & \TODO{изм.} & \TODO{n} \\
\bottomrule
\end{tabular}
\end{table}

\TODO{редовете на таблицата се добавят, след като се определи наборът от заявки в
Раздел V; всяка заявка получава ред}

\subsection{Цена на извеждането по RDFS}

Двата режима на слоя за извеждане се сравняват помежду си и с изпълнение без
извеждане. Очакваният вид на резултата е компромис между еднократна и повтаряща се
цена, но посоката и големината му се записват едва след измерване.

\TODO{таблица за сравнение на двата режима на извеждане}

\subsection{Анализ}

\TODO{анализ на получените резултати: къде избраните проектни решения се потвърждават,
къде не, и коя от тях е причината за наблюдаваната разлика}
